\chapter{Symmetric powers over a commutative ring}\label{ch:la01-symmetric}

Throughout the chapter \(R\) is a commutative ring with identity, all
algebra maps preserve the identity, and \(M\) is an \(R\)-module.
The construction uses the balanced tensor product and its
associativity from \emph{Algebra Core}
\cite[ALG-03-C02-S01]{p103}. The field case, including the monomial
basis, is developed in \emph{Linear Algebra}
\cite[entries 15.symmetric and 15.symmetric\_basis]{d05}. Here the
coefficient ring and the module may be arbitrary.

\section{The symmetric algebra and its universal properties}

\begin{definition}[LA01-C05.01. The symmetric algebra]
\label{def:la01-symmetric-algebra}
The tensor algebra is \(T_R(M)=\bigoplus_{n\ge 0}M^{\otimes_R n}\),
with \(M^{\otimes_R 0}=R\) and product given by concatenation of
tensors. The symmetric algebra \(\operatorname{Sym}_R(M)\) is its
quotient by the two-sided ideal \(J\) generated by
\(x\otimes y-y\otimes x\) for \(x,y\in M\). The degree-\(n\) component
is the symmetric power \(\operatorname{Sym}^n_R(M)\).
\end{definition}

The ideal \(J\) is homogeneous because its generators have degree
two. More explicitly, a homogeneous piece of a product
\(a(x\otimes y-y\otimes x)b\) is a sum of such products where \(a,b\)
are homogeneous and their degrees add to the required degree minus
two. There are no relations in degrees zero and one. Thus
\(\operatorname{Sym}^0_R(M)=R\) and
\(\operatorname{Sym}^1_R(M)=M\) canonically. We write \(x_1\cdots x_n\)
for the image of a pure tensor, with empty product \(1\). Scalars
commute with everything, and degree-one generators commute by the
defining relations. Since these generate the quotient as an
\(R\)-algebra, it is commutative. See \cite[Algebra, Section 10.13]{stacks-sym}.

\begin{proposition}[LA01-C05.02. Universal symmetric multilinear maps]
\label{prop:la01-symmetric-multilinear}
For every \(n\ge 0\) and every \(R\)-module \(N\), the assignment
\(g\mapsto ((x_1,\ldots,x_n)\mapsto g(x_1\cdots x_n))\) gives a natural
bijection from \(R\)-linear maps \(\operatorname{Sym}^n_R(M)\to N\)
to symmetric \(R\)-multilinear maps \(M^n\to N\). For \(n=0\) the
right-hand side means the choices of a single element of \(N\).
\end{proposition}
\begin{proof}
For \(n\ge 1\), iteration of the tensor universal property
gives a unique linear map on \(M^{\otimes_R n}\) extending any
multilinear map on \(M^n\). In degree \(n\ge 2\), the homogeneous
description of \(J\) shows that \(J\cap M^{\otimes_R n}\) is spanned
by differences of pure tensors related by one adjacent swap.
Indeed the factors before and after the generator can be taken
to be pure tensors, which span their respective tensor powers.
A linear map kills these differences precisely when its associated
multilinear map is invariant under adjacent swaps. Those swaps
generate the whole permutation group. This proves the bijection
for \(n\ge 2\), and for \(n=1\) both descriptions give all linear
maps \(M\to N\). For \(n=0\) use \(g\mapsto g(1)\) on
\(\operatorname{Hom}_R(R,N)\).

The construction is natural in \(N\) by postcomposition. A
linear map \(M\to M'\) induces tensor powers that take adjacent
differences to adjacent differences. The induced maps on the
quotients give naturality in \(M\) by precomposition.
\end{proof}

\begin{proposition}[LA01-C05.03. The free commutative algebra property]
\label{prop:la01-symmetric-free}
For any commutative \(R\)-algebra \(A\), each \(R\)-linear map
\(f\colon M\to A\) extends uniquely to an \(R\)-algebra map
\(\operatorname{Sym}_R(M)\to A\).
\end{proposition}
\begin{proof}
For every \(n\), the product \(f(x_1)\cdots f(x_n)\) is symmetric
and multilinear. Proposition~\ref{prop:la01-symmetric-multilinear}
extends it to a linear map on the degree-\(n\) component.
For degree zero this gives the structure map \(R\to A\).
The direct sum of these maps respects multiplication on products
of pure classes, where it agrees with multiplication of the images
of all their factors in \(A\). By bilinearity it respects
multiplication on the whole algebra, and its degree-zero component
takes the unit to the unit. The images of \(R\) and of \(M\)
determine the map on all products and linear combinations of
generators, giving uniqueness.
\end{proof}

\begin{theorem}[LA01-C05.04. The monomial basis of a finite free module]
\label{thm:la01-symmetric-basis}
Suppose \(M\) is free with basis \(e_1,\ldots,e_r\), allowing \(r=0\).
The assignment \(e_i\mapsto X_i\) gives a graded \(R\)-algebra
isomorphism
\[
  \operatorname{Sym}_R(M)\cong R[X_1,\ldots,X_r],
\]
where every \(X_i\) has degree one. For every \(n\ge 0\), the
elements \(e_1^{a_1}\cdots e_r^{a_r}\), with all \(a_i\ge 0\) and
\(\sum_i a_i=n\), are an \(R\)-basis of \(\operatorname{Sym}^n_R(M)\).
When \(r\ge 1\) there are \(\binom{n+r-1}{n}\) basis indices. For
\(r=0\), the basis consists of \(1\) in degree zero and is empty
in positive degrees.
\end{theorem}
\begin{proof}
Send the basis linearly to the variables and extend using
Proposition~\ref{prop:la01-symmetric-free}. Sending \(X_i\) to \(e_i\)
defines the reverse map by the universal property of the polynomial
algebra. Each composition fixes the scalars and all the algebra
generators on its side, so both compositions are identities.
The maps preserve degree. The formal polynomial algebra is the free
\(R\)-module on its formal monomials, and its degree pieces are
spanned by the monomials of the indicated total degree. For \(r\ge 1\),
the exponents can be recorded by placing \(r-1\) separators among
\(n+r-1\) positions, with the \(n\) other positions recording the
degree. For \(r=0\) the polynomial algebra is \(R\) in degree zero.
This gives the asserted bases in every degree.
\end{proof}

The basis statement over a commutative ring also follows from the
multilinear approach in \cite[Theorem 1.2]{conrad-sym}.

\section{Extension of scalars}

\begin{theorem}[LA01-C05.05. Arbitrary base change]
\label{thm:la01-symmetric-base-change}
Let \(R\to S\) be any homomorphism of commutative rings with
identity and let \(M_S=S\otimes_R M\). The graded \(S\)-algebra map
\[
  S\otimes_R \operatorname{Sym}_R(M)
  \longrightarrow \operatorname{Sym}_S(M_S),
  \qquad
  s\otimes (x_1\cdots x_n)\longmapsto
       s(1\otimes x_1)\cdots(1\otimes x_n)
\]
is an isomorphism and gives a comparison in every degree \(n\ge 0\).
It is natural in \(M\) and compatible with compositions of
ring maps under the usual associativity of scalar extension.
\end{theorem}
\begin{proof}
The \(R\)-linear map \(x\mapsto 1\otimes x\), with values in
degree one of \(\operatorname{Sym}_S(M_S)\), extends to an
\(R\)-algebra map by Proposition~\ref{prop:la01-symmetric-free}.
Multiplying its values by \(s\) defines the indicated map on
simple tensors, by the balanced tensor universal property. This
map respects scalar multiplication and algebra multiplication.
In the opposite direction, \(s\otimes x\mapsto s\otimes x\) with
values in degree one of \(S\otimes_R\operatorname{Sym}_R(M)\)
is the scalar extension of the degree-one inclusion. It extends
to an \(S\)-algebra map by the same proposition over \(S\).
The \(1\otimes x\) generate both algebras over \(S\), and the
two compositions fix these elements and the scalars. The maps
are therefore inverse. They preserve degree, and tensor products
commute with direct sums, so the degreewise comparisons follow.

For a module map \(f\colon M\to M'\), the two routes send the
generator \(1\otimes x\) to \(1\otimes f(x)\). For \(S\to U\),
the two routes after scalar extension send \(1\otimes(1\otimes
x)\) to the degree-one generator \(1\otimes x\) over \(U\).
In both cases agreement on the generators proves the claim.
\end{proof}

The base-change comparison is the module version of the pullback
formula for symmetric algebras
\cite[Modules, Lemma 17.21.3]{stacks-sym}. For exterior powers,
\emph{Advanced Commutative Algebra} gives the alternating quotient
construction and its coefficient comparison in RK-03 and RK-04
\cite[RK-03 and RK-04]{commalg}. In degree \(q\) the map is
\[
  s\otimes(x_1\wedge\cdots\wedge x_q)
  \longmapsto s(1\otimes x_1)\wedge\cdots\wedge(1\otimes x_q).
\]
The alternating condition is vanishing on repeated inputs,
including in characteristic two. RK-04 applies the comparison to
the terms of the Koszul complex for every ring map and uses
flatness in the statement about Koszul homology.

\begin{example}[LA01-C05.06. A torsion module and nonflat base change]
\label{ex:la01-symmetric-torsion}
Take \(R=\mathbb Z\), \(M=(\mathbb Z/2\mathbb Z)e\) with generator
\(e\), and \(S=\mathbb F_2\). A linear map \(M\to A\) into a
commutative \(\mathbb Z\)-algebra is specified by an element of
\(A\) annihilated by \(2\). The same choices specify
\(\mathbb Z\)-algebra maps \(\mathbb Z[X]/(2X)\to A\) via \(X\).
The canonical maps induced by \(e\leftrightarrow X\) are inverse
on the generators, so
\(\operatorname{Sym}_{\mathbb Z}(M)\cong\mathbb Z[X]/(2X)\).
Its degree-zero part is \(\mathbb Z\), and each positive degree
is generated by \(X^n\) with coefficient module \(\mathbb Z/2\mathbb Z\).
After scalar extension the quotient becomes \(\mathbb F_2[X]\),
in agreement with the rank-one description of
\(\operatorname{Sym}_{\mathbb F_2}(M_S)\).
This extension is nonflat because tensoring
\(\mathbb Z\xrightarrow{\,2\,}\mathbb Z\) with \(\mathbb F_2\)
sends an injection to the zero map on a nonzero module.
\end{example}

\begin{example}[LA01-C05.07. The quotient in characteristic two]
\label{ex:la01-symmetric-char-two}
Let \(M=\mathbb F_2 e_1\oplus\mathbb F_2 e_2\) and let \(\tau\)
exchange the factors of \(M^{\otimes 2}\) over \(\mathbb F_2\).
In the tensor basis the fixed subspace has basis
\[
  e_1\otimes e_1,\qquad e_2\otimes e_2,\qquad
  e_1\otimes e_2+e_2\otimes e_1.
\]
Indeed the coefficients of the two mixed tensors must be equal.
The quotient map takes the third basis element to
\(2e_1e_2=0\), even though \(e_1e_2\) is a nonzero basis monomial
of the quotient. Thus the map from the fixed tensors to the
symmetric quotient is neither injective nor surjective. The
averaging construction in \cite[entry 15.normalization]{d05}
uses invertible factorials in its field setting, while the
quotient and the monomial basis above hold in every characteristic.
\end{example}
